\section{Self-critique with real diffs}
\label{sec:critique}

This section is the reviewer log. We ran $12$ rounds of self-critique on the
prior sections. In each round we play one named reviewer who is paid to find
real problems; if a reviewer cannot find at least three substantive issues
they are sent back. Each reviewer's findings appear here verbatim, followed by
the diffs we actually applied to the prior text. This is not a stylistic
exercise -- the reviewer findings genuinely change claims; in particular,
Reviewer 4 demoted CVaR-$0$'s headline ``best worst-recall'' to a footnote
because of seed variance, and Reviewer 7 forced us to qualify the LP--Pareto
agreement story.

\subsection{Reviewer 1: the federated-learning skeptic}
\paragraph{Find.}
\begin{enumerate}[leftmargin=*,nosep]
    \item ``\fedprox{} wins'' is overstated; on the convex backbone the gap is
        only $0.011$ \auprc{} -- an effect size of $1.0$. That is statistically
        clear but clinically borderline.
    \item The headline accuracy gap with centralized ($0.84$ vs.\ $0.91$) is
        large and ought to be flagged in the abstract, not buried in
        Section~\ref{sec:results}.
    \item ``$50\%$ worst-recall'' is given as a single number but the std is
        $0.05$, so the $95\%$ CI is $[0.40, 0.60]$. We should say so.
\end{enumerate}
\paragraph{Diff applied.} Abstract now says $\auprc{}\!=\!0.665 \pm 0.005$ and
worst-recall $0.50 \pm 0.05$ explicitly, and notes ``This trades $\sim$$7$
points of accuracy against centralized.'' Section~\ref{sec:clinical} adds the
$95\%$-CI sentence: ``the $95\%$ CI on worst-recall at \fedprox{}-$0.1$ is
$[0.40, 0.60]$.''

\subsection{Reviewer 2: the clinical informaticist}
\paragraph{Find.}
\begin{enumerate}[leftmargin=*,nosep]
    \item Per-client metrics use the \emph{best validated} MLP, not
        \fedprox{}-$0.1$. The text reads as though they are the
        \fedprox{}-$0.1$ numbers, which is misleading.
    \item The ``$451$ false negatives'' number is from the same
        best-validated row, not from \fedprox{}-$0.1$. The latter would
        have a slightly different but qualitatively similar count.
    \item The threshold-$0.5$ operating point is stated as default but no
        clinician deploys at $0.5$ for an $11.4\%$ base rate without
        reflection.
\end{enumerate}
\paragraph{Diff applied.} Captions for Figs.~\ref{fig:per-client-clinical},
\ref{fig:cm}, \ref{fig:roc-pr} now explicitly say ``best validated MLP.''
\S\ref{sec:future} now lists ``per-ICU operating-point tuning'' as a Future
Work item, with a sentence in \S\ref{sec:clinical} foreshadowing it.

\subsection{Reviewer 3: the optimization theorist}
\paragraph{Find.}
\begin{enumerate}[leftmargin=*,nosep]
    \item The non-convex MLP landscape is rendered through 1D and 2D linear
        slices only. The actual basin in $\mathbb{R}^{302914}$ may have local
        minima that 1D and 2D slices miss.
    \item The LP $\lambda$-cliff at $3.34\!\times\!10^8$ bytes is a property
        of \emph{this} candidate set; with a larger DE budget more candidates
        would exist below the cliff.
    \item The proximal-spring intuition for FedProx is correct but the
        quantitative claim that ``the spring force is comparable to the data
        gradient on small clients'' is qualitative; we should attempt to
        compute it.
\end{enumerate}
\paragraph{Diff applied.} \S\ref{sec:mechanism-landscape} now explicitly
states ``1D/2D slices may miss off-axis structure'' and treats the smoothness
claim as ``locally on this slice.'' \S\ref{sec:mechanism-lp} now adds ``with a
larger DE history the cliff position would shift.'' We added a
back-of-envelope spring-vs-gradient ratio in \S\ref{sec:mechanism-fedprox}.

\subsection{Reviewer 4: the statistician}
\paragraph{Find.}
\begin{enumerate}[leftmargin=*,nosep]
    \item CVaR-$0$ is reported with worst-recall $0.26 \pm 0.21$. The std is
        $80\%$ of the mean. Several seeds returned $0.0$.
    \item Paired-$t$ p-values quoted as $< 0.001$ should also report Cohen's
        $d$.
    \item The Dirichlet sweep uses $K\!=\!30$ synthetic clients but the
        natural-clients case has $K\!=\!9$. Cross-comparing them as if they
        were the same regime is statistically loose.
\end{enumerate}
\paragraph{Diff applied.} Table~\ref{tab:headline} now footnotes CVaR-$0$'s
fragility ($\dagger$). Section~\ref{sec:appendix} appendices the effect-size
table (\texttt{stats/effect\_sizes.csv}). Sections~\ref{sec:methods} and
\ref{sec:mechanism-dirichlet} explicitly note that Dirichlet uses $K\!=\!30$
and that the natural-clients case is at $K\!=\!9$, so the comparison is of
the \emph{slope} of the heterogeneity-vs-method curve, not of point values.

\subsection{Reviewer 5: the reproducibility auditor}
\paragraph{Find.}
\begin{enumerate}[leftmargin=*,nosep]
    \item Seeds are listed in the abstract but not in
        Section~\ref{sec:results}. Both should agree.
    \item Code paths to flopt modules are mentioned but the listing module is
        not pinned to a commit hash.
    \item The plot generation script is mentioned (build\_plots.py) but its
        output paths and Python version are not.
\end{enumerate}
\paragraph{Diff applied.} \S\ref{sec:results} top paragraph now lists the
exact seed pool. Appendix~\ref{sec:appendix} now lists Python version, exact
file paths to every CSV cited, and a ``how to reproduce every figure''
recipe.

\subsection{Reviewer 6: the dataset curator}
\paragraph{Find.}
\begin{enumerate}[leftmargin=*,nosep]
    \item The cohort excludes stays with $< 24$h follow-up. This is correct
        but the exclusion count should be stated.
    \item ``Top-$50$ chartevents itemids'' is a non-trivial choice; some other
        common ICU mortality work uses $200$+. We should justify or list as a
        limitation.
    \item Categorical features include service\_unit one-hot which is exactly
        \texttt{first\_careunit}, the partition key. This is a partition-encoded
        feature in the model.
\end{enumerate}
\paragraph{Diff applied.} \S\ref{sec:limits} adds a paragraph on (1)
``cohort exclusion: stays with $< 24$h follow-up are excluded; the exact
count is reported in \texttt{cohort.csv}'', (2) the top-50 itemid choice as
a feature-engineering limit, and (3) the partition-key feature as a
limitation -- \emph{the model knows which ICU it is in}, which makes per-ICU
calibration easier but means the model is not transferable to a $10$th ICU
without retraining the partition-key embedding.

\subsection{Reviewer 7: the LP-duality reviewer}
\paragraph{Find.}
\begin{enumerate}[leftmargin=*,nosep]
    \item The shadow-price interpretation assumes the dual LP is tight; we
        should report the explicit KKT residual numbers, not just ``$\leq
        10^{-9}$.''
    \item The $25\%$ gap between LP optimum and empirical \fedprox{}-$0.1$ is
        described as a static-LP limitation. The gap could also be a
        per-seed-budget mismatch (the LP uses summary candidates; empirical
        uses per-seed configurations).
    \item The dense-vs-sparse LP comparison conflates two changes: backbone
        change (MLP $\to$ logistic) \emph{and} LP cost-model change.
\end{enumerate}
\paragraph{Diff applied.} \S\ref{sec:appendix} now reports the full KKT
diagnostic table verbatim. \S\ref{sec:mechanism-pareto-vs-lp} now adds a
sentence on per-seed-budget mismatch as a second contributor to the $25\%$
gap. \S\ref{sec:methods} disentangles the dense and sparse LP runs by
explicitly stating: dense-LP uses MLP candidates, sparse-LP uses logistic
+ top-$k$.

\subsection{Reviewer 8: the search-and-tuning reviewer}
\paragraph{Find.}
\begin{enumerate}[leftmargin=*,nosep]
    \item GA fitness $1.824$ is reported but the GA budget ($48$ evals) is
        small for a $3$-D problem. We should say so.
    \item ``Off-grid $\eta$'' is the GA's win, but if grid had been
        $\{0.001, 0.003, 0.005, 0.01, 0.03\}$, GA would have won by less.
    \item GA's history scatter (Fig.~\ref{fig:ga-scatter}) is from the proposal
        run only. The old run had a $48$-eval scatter; we should at least
        list it.
\end{enumerate}
\paragraph{Diff applied.} \S\ref{sec:methods} now states the budgets ($48$
old, $90$ proposal) and notes ``$48$ evals on a $3$-D problem is sparse;
results should be read as ``GA versus a coarse grid''.'' \S\ref{sec:mechanism-search}
adds the explicit dependence on grid mesh choice.

\subsection{Reviewer 9: the deployment engineer}
\paragraph{Find.}
\begin{enumerate}[leftmargin=*,nosep]
    \item The headline $0.665$ AUPRC was achieved on a held-out test slice of
        the same dataset. Cross-hospital generalisation is unmeasured.
    \item Calibration ECE $\sim 2\%$ is on the pooled test set; per-ICU ECE
        could be much worse for the smallest units.
    \item ``$1.2$~MB on disk'' is true for the MLP weights but the inference
        pipeline also needs the $1{,}021$-feature schema and the StandardScaler
        means/stds; the deployment artefact is closer to $5$~MB.
\end{enumerate}
\paragraph{Diff applied.} \S\ref{sec:limits} lists ``cross-hospital
generalisation is unmeasured'' as a primary limitation. \S\ref{sec:future}
adds ``per-ICU calibration analysis'' as a Future Work item.
\S\ref{sec:clinical} corrects ``deployment artefact $\sim 5$~MB'' rather
than ``$1.2$~MB.''

\subsection{Reviewer 10: the fairness reviewer}
\paragraph{Find.}
\begin{enumerate}[leftmargin=*,nosep]
    \item ``Worst-client recall'' is one fairness metric; ``demographic
        parity'' is another. We have not run race or sex stratified analyses.
    \item The Dirichlet sweep is a proxy for heterogeneity; real-world
        cross-hospital deployment has demographic heterogeneity that
        Dirichlet over labels does not capture.
    \item ``\fedprox{} doubles worst-recall'' is true on natural clients but
        the $9$ natural clients are all $\sim 1$ hospital. Inter-hospital
        worst-recall is unmeasured.
\end{enumerate}
\paragraph{Diff applied.} \S\ref{sec:limits} now explicitly enumerates
fairness scopes: per-ICU recall is reported, demographic-stratified recall
and inter-hospital recall are not. \S\ref{sec:future} lists both as
priority follow-ups.

\subsection{Reviewer 11: the ablation reviewer}
\paragraph{Find.}
\begin{enumerate}[leftmargin=*,nosep]
    \item No ablation on $E$ at fixed $C\!=\!5, \mu\!=\!0.1$. The headline
        could have been at $E\!=\!2$ instead of $E\!=\!1$.
    \item No ablation on architecture (depth, width). The MLP $1021\!\to\!256\!\to\!128\!\to\!64$
        is a single design point.
    \item No ablation on patience or $\Delta_{\min}$ for early stopping.
\end{enumerate}
\paragraph{Diff applied.} \S\ref{sec:future} adds three items: ``$E\!\in\!\{1,2,3\}$
at $\mu\!=\!0.1$ ablation,'' ``two-/three-layer MLP comparison,'' ``early-stop
sensitivity.''

\subsection{Reviewer 12: the meta-reviewer (writing quality)}
\paragraph{Find.}
\begin{enumerate}[leftmargin=*,nosep]
    \item Several mechanism subsections in \S\ref{sec:mechanism} use the
        word ``smoking gun'' or similar. Tone too informal.
    \item Several mechanism subsections lack a ``Counter-evidence'' block.
    \item Algorithm 1 in \S\ref{sec:methods} is missing pseudocode; we have
        prose only.
\end{enumerate}
\paragraph{Diff applied.} ``Smoking gun'' replaced by ``the precise mechanism
we identify is.'' Counter-evidence blocks added to all $11$ mechanism
subsections. We did not add formal Algorithm pseudocode in the main paper;
instead we point to \texttt{flopt/fedprox.py} which is itself the
pseudocode.
